\section{CL-TAV Method}
\label{sec:method}

\subsection{Notation and assumptions}
S0--S10 are logical duties, not line numbers. The decision index \(k\), the
billed charge count and the retry counter are distinct. SessionContext
\(\Gamma\) is the unique session state. HistoryHandle \(\eta\) carries
whole-history compatibility. The select-time snapshot \(\xi\) freezes
\texttt{actionId}, \texttt{actionKind}, \(q_{\mathrm{used}}\), history version,
classes used at select, and uncertainty reference \(U\). A later confirmed
successor summary is a different field. Effect classes are
\texttt{CONFIRMED-NOT-SENT} and \texttt{UNKNOWN-EFFECT}. PredictionGapError is
TEST-scoped.

\subsection{Overall closed loop}
Algorithm~\ref{alg:cltav-01} is a typed abstraction of ALG-CLTAV-01. Wrappers
\textsc{PredictCurrent}, \textsc{InterpretOutcome} and
\textsc{ResolveOutcome} are ALG-CLTAV-02/03/04 in the supplement. S9 calls
\textsc{CommitCompatibleUpdate} (ALG-CLTAV-07), with
\(\chi=(\Gamma,\eta,H)\), which validates the history
version, applies \IFhist{} from \(\xi\), returns
\((\textit{status},\eta,H,\Gamma)\) and commits a confirmed successor summary.
A bare \IFhist{} call is not that wrapper. Resolution is branched on
\(\delta.\mathrm{control}\) after adopting \(\delta.\Gamma\), \(\delta.\eta\)
and \(\delta.\mathrm{outcome}\).

% Compact display of ALG-CLTAV-01. Same duties and interfaces as the
% authoritative module; not a second control algorithm.
\begin{algorithm}[t]
\caption{CL-TAV top-level session (paper Alg.~1 = ALG-CLTAV-01)}
\label{alg:cltav-01}
\SetKwInput{KwIn}{Require}
\SetKwInput{KwOut}{Ensure}
\KwIn{\(H_0=\{h_{\mathrm{normal}}\}\cup H_{\mathrm{single}}\); \(\Gamma\); \(\eta\); one mode BUDGET\((c_{\min},B)\) xor ROUNDS\((K_{\max})\); library \(L\); uncertainty \(U\)}
\KwOut{stop class, final \(H\), \(\eta\), trace; singleton \(h_{\mathrm{normal}}\) is not protocol PASS}
\textsc{Initialize}; \While{true}{
  \lIf{\IFres{} already decided}{\Return that stop}
  \(\textit{pred}\leftarrow\textsc{PredictCurrent}\) via \IFpred\;
  \(\textit{dec}\leftarrow\textsc{SelectAndAdmit}\) via \IFsel\;
  \lIf{\(\textit{dec}\) is EXIT or SPEC-ERROR}{\Return its reason}
  freeze snapshot \(\xi\) (action, \(q_{\mathrm{used}}\), history version, classes)\;
  charge \(\xi\) once; \(\textit{er}\leftarrow\) \IFexec\((\Gamma,\xi)\)\;
  \(z\leftarrow\textsc{InterpretOutcome}\) via \IFobs{} and ALG-CLTAV-06\;
  \(\delta\leftarrow\textsc{ResolveOutcome}\); adopt returned \((\Gamma,\eta,z)\)\;
  \lIf{\(\delta=\texttt{STOP}\)}{\Return reason}
  \lIf{\(\delta=\texttt{RETRY}\)}{\textbf{continue}}
  \If{\(z\) has valid \(I_z\) or a confirmed successor summary}{
    commit via \IFhist{} from \(\xi\), not from the post-effect \(q\)\;
  }
  \lIf{a P-stop is set}{\Return it}
}
\end{algorithm}


\subsection{Interpretation, compatibility and selection}
Prediction and interpretation share \(U\). Let \(O(h,t,q)\) be the conservative
outer set of feasible observations of hypothesis \(h\), and let \(Z_k\) be the
measured feasible-observation set after timing, ownership and pairing checks.
The interpretation/whole-history bridge induces the compatible hypothesis
class \(C_k\) (the accepted interface field \(z.I_z\)); \(C_k\) is not a
measurement set. A currently valid class is the nonempty projection
\(H\cap C_k\). Compatible update on a valid class is
\begin{equation}
\label{eq:update}
H_{k+1}=H_k\cap C_k,\qquad H_{k+1}\subseteq H_k.
\end{equation}
\texttt{ERROR} and identity Prep/Recover leave \(H\) unchanged. The one-step
score on currently valid nonempty classes of an affordable distinguishing
TEST \(t\) is
\begin{equation}
\label{eq:score}
s(t)=\max_{C\in\mathcal{C}_t}\lvert C\rvert.
\end{equation}
Empty \(S\) is never \(\arg\min\). Timing uses T5~\cite{jcgm106}:
\(I_{\mathrm{obs}}\subseteq I_r\) is PASS; \(I_{\mathrm{obs}}\cap I_r=\varnothing\)
is FAIL; otherwise INCONCLUSIVE. An ambiguous pair is \texttt{ERROR}, not IUT
FAIL.

\subsection{Bounded reference-kernel contract}
The first implementation target is a finite explicit-path kernel, not a
general solver. Each frontier entry contains a stable hypothesis ID, control
state, canonical exact-rational difference constraints, finite typed store, path length,
whole-history provenance, and a known/conservative-unknown status. It accepts
only TRUE, conjunction, state equality, typed-field equality and rational
interval guards, with simultaneous finite assignment and exact zero-clock-reset
updates. OR, NOT and calls are outside this first domain and are rejected as
SPEC-ERROR before expansion; no normalization is left to an implementer.
The sole normative numeric form is a numerator and strictly positive
denominator pair; interval bounds carry explicit open/closed flags, and any
fraction printed in a review view is display-only.  A typed adapter either
maps the supported M2 subset into this AST or rejects the model before
expansion, so M2 expressions are not a second authoritative encoding.
Propagation validates the select-time snapshot, expands enabled transitions,
intersects observation and timing constraints, and retains every feasible
successor. Entries merge only when hypothesis, control state, clock
constraint, store and history provenance are identical.

Configured bounds cover hypotheses, actions, observation classes, path length
and per-hypothesis frontier size. Unsupported syntax or an exhausted path or
frontier bound returns the interface's named GAP, SPEC-ERROR, or conservative
unknown result without deleting a hypothesis. Only prediction, selection,
history update, equivalence and stop selection consume kernel results;
execute, interpretation and Prep/Recover do not. Feasibility alone never
establishes equivalence. Unsupported syntax, computation unknown and limits
preserve \(H\), the summary and charge state, irrespective of a later stop label.
A valid successor equals the registered compatibility intersection and is a
subset of \(H\). Its empty result forces Stop-Empty and its singleton result
forces Stop-Singleton before 645, equivalence or budget exits. Stop-Empty thus
requires an empty history proved by complete in-domain intersection. This contract maps every
reachable result and marks every other interface/result pair unreachable; it does not claim
completeness, polynomial complexity, or deployed code.

This contract is intentionally weaker than a universal uncertain-observation
diagnosis theorem. Existing work already supplies possible-run semantics and
an exact incremental passive diagnoser for bounded latency, drift/offset and
out-of-order observation under exact-once and timed-automaton assumptions
\cite{koehl2023}, as well as exact monitoring/testing for a different bounded
communication-delay model~\cite{fraenzle2026}. CL-TAV claims only the consistency
and retention property below under R1--R5. Detected capture inadequacy or
invalid evidence, ambiguous ownership or pairing, unsupported syntax and
exhausted bounds retain their existing named uncertain, error or resource
outcomes and do not establish an IUT failure. This does not claim that every
missing, duplicate or spurious event is detectable: undetected out-of-domain
pollution may remain compatible with an in-domain candidate. Conversely, with
qualified evidence and a valid still-active obligation, expiry at the declared
timing boundary retains the existing no-response FAIL; a legal observation
class whose complete in-domain intersection is empty retains Stop-Empty. No
source theorem is transplanted.

\subsection{Conditional properties}
\textbf{Monotonicity.} Equation~\eqref{eq:update} gives
\(H_{k+1}\subseteq H_k\). Without a valid class there is no add or delete. An
in-domain singleton is not protocol PASS.

\textbf{Retention.} Assume (R1) \(h^\star\in H_0\); (R2) one state path is
consistent with the executed history; (R3) \(O(h^\star,t_k,q_k)\) is a
conservative outer set; (R4) the true observation belongs to both
\(O(h^\star,t_k,q_k)\) and \(Z_k\); and (R5) ownership and pairing are valid.
The declared interpretation and whole-history compatibility contract therefore
places \(h^\star\) in \(C_k\). Thus \(h^\star\in H_k\cap C_k\), and induction
retains it in \(H_{k+1}\). Out-of-domain faults may still mimic an in-domain
survivor.

\textbf{One-step minimax.} On nonempty finite \(T_{\mathrm{cur}}\), \(t^\star\)
is a least element of \((s(t),\mathrm{cost}(t),\mathrm{id})\). Diagnostic-state
feedback and timed active control already exist~\cite{hu2020,miao2025}; this
statement concerns only the declared current choice. It is not a global policy,
does not prove future distinguishability, and does not guarantee shrinkage for
every observation class.

\textbf{Termination.} Every charged TEST, Prep, Recover and \texttt{ERROR}
spends the selected mode, hence at most \(\lfloor B/c_{\min}\rfloor\) or
\(K_{\max}\) billed executions. Wall-clock waits still need their own timeouts;
those timeouts are not claimed implemented.

Exclusive named stops are Empty, Singleton, 645, Equivalent, Budget,
Stop-Error and Stop-NoDistinguisher. Stop-Error is not Stop-Budget.
Admit reasons A2--A5 and post-update P1--P5 are tabulated in the supplement.
CONFIRMED-NOT-SENT does not assign UNKNOWN. UNKNOWN-EFFECT must not be read
back as the previous KNOWN summary. Unconfirmed Recover with \(R_{\max}=1\)
and remaining budget is Stop-Error. The update disposition is evaluated before
the mutable stop label: a rejected/error update preserves \(H\), while a valid
empty or singleton update forces Stop-Empty or Stop-Singleton, respectively.
Kernels remain unimplemented interfaces.
